% =============================================================================
% 04_RESULTS — Thematic Synthesis Findings
% =============================================================================
% Page target: 2–4 pages
% Structure: 3 themes with sub-themes — synthesis, not catalogue
% =============================================================================

\section{Results}
\label{sec:results}
\addcontentsline{toc}{section}{Results}

% =============================================================================
% OVERVIEW
% =============================================================================

From an initial corpus of 246 publications identified through systematic
search across PubMed, Semantic Scholar, arXiv, and Google Scholar, [N]
papers were retained after deduplication, abstract screening, and full-text
assessment against inclusion and exclusion criteria. The reviewed literature
reveals three overarching findings at the intersection of health AI
evaluation and governance: a pervasive \textbf{ecological validity gap}
in evaluation methodology; a structural \textbf{governance capacity deficit}
in health systems; and a systematic \textbf{pilot-to-routine transition
failure} rooted in organizational rather than technical barriers. These three
findings are analytically interdependent rather than independent — the
ecological validity gap generates misleading validation evidence that reduces
the urgency of governance capacity-building; governance frameworks built on
that evidence are structurally misaligned with real-world complexity; and
that misalignment produces the pilot trap as a downstream consequence.
than technical barriers.

% =============================================================================
% THEME 1: EVALUATION METHODS — Ecological Validity Gap
% =============================================================================

\subsection{Evaluation Methods: The Ecological Validity Gap}
\label{sec:results_evaluation}

% -----------------------------------------------------------------------------
% 1.1 Single-site, retrospective validation dominates
% -----------------------------------------------------------------------------

\subsubsection{Limited External Validation and Retrospective Bias}
\label{sec:eval_retrospective}

A consistent methodological limitation across the reviewed corpus is the
predominance of single-site, retrospective evaluation designs. Studies
evaluating deep learning models for medical imaging — including models for
colorectal cancer microsatellite instability prediction, diabetic retinopathy
screening, and chest X-ray interpretation — routinely report Area Under the
Curve (AUC) values exceeding 0.85 in internal validation cohorts but show
substantial performance decay when tested on external datasets from different
institutions or populations \citep{pham2025ethical}. This pattern — high
internal validity, limited external generalizability — reflects a structural
evaluation problem rather than a model-specific failure.

The CT-based MSI prediction study exemplifies this gap with quantitative
precision: the authors report AUCs of 0.882 in their training cohort,
declining to 0.768 in the internal validation cohort, 0.803 in external
validation cohort 1, and 0.751 in external validation cohort 2 — a
progressive performance degradation across increasingly distant populations
that illustrates the generalizability challenge \citep{pham2025ethical}.
Importantly, even when external validation is attempted, it typically
relies on retrospective data from the same health system, capturing little
of the population diversity, data quality variation, and workflow integration
complexity of real clinical environments.

% -----------------------------------------------------------------------------
% 1.2 Metrics misaligned with clinical priorities
% -----------------------------------------------------------------------------

\subsubsection{Metric Design and Clinical Relevance}
\label{sec:eval_metrics}

Standard AI evaluation metrics — accuracy, sensitivity, specificity, AUC —
are poorly calibrated to the clinical context in which AI systems operate.
\citet{turchi2025ecological} demonstrate this with particular clarity in their
study of AI-assisted oral cancer diagnosis: experienced dental specialists
identifying real-world clinical diagnosis as requiring contextual information
(beyond static images), longitudinal reasoning (beyond single-session tasks),
and harm-asymmetric evaluation (a false-negative in cancer diagnosis has a
different clinical weight than a false-positive). Standard AI benchmarks
collapse these distinctions into a single accuracy number.

The TRAC-Pain protocol \citep{bailes2026trac} offers a methodological
counterpoint. By combining continuous wearable sensor data (physiological,
sleep, and physical activity metrics) with ecological momentary assessments
(EMAs) administered multiple times daily over a 12-week period, the study
designs evaluation metrics that reflect the temporal dynamics of the pain
experience — including pain interference, mood, and fatigue — rather than
relying on point-of-care assessments alone. This represents an important
step toward ecologically grounded evaluation, though prospective longitudinal
designs of this kind remain rare in the literature.

% -----------------------------------------------------------------------------
% 1.3 Continuous monitoring and drift detection
% -----------------------------------------------------------------------------

\subsubsection{Algorithmic Drift and Post-Market Monitoring}
\label{sec:eval_drift}

A critical and understudied dimension of AI evaluation is algorithmic drift:
the phenomenon by which AI systems that were valid at the time of deployment
gradually or abruptly degrade in performance as input data distributions
shift — due to changes in patient populations, healthcare practices, laboratory
equipment, or the AI system's own continuous learning updates. The FDA
perspective paper \citep{warraich2024fda} identifies continuous learning AI
systems as requiring special regulatory attention precisely because the
traditional pre-market validation paradigm cannot guarantee post-deployment
performance. Yet mandatory post-market performance monitoring with standardized
drift detection protocols is conspicuously absent from most regulatory
frameworks reviewed.

\citet{abulibdeh2025illusion} argue that this absence constitutes a
\textbf{systemic safety gap}: patients and clinicians are provided no
mechanism to know whether an AI system that was validated two years ago
remains fit for clinical purpose today. The authors call for mandatory
real-time monitoring requirements as a condition of market authorization,
with clear escalation protocols and performance thresholds that trigger
mandatory re-validation or market withdrawal.

% =============================================================================
% THEME 2: GOVERNANCE STRUCTURES — Capacity Deficit and Adaptive Regulation
% =============================================================================

\subsection{Governance Structures: The Capacity Deficit}
\label{sec:results_governance}

% -----------------------------------------------------------------------------
% 2.1 Pre-market models assume pre-existing governance capacity
% -----------------------------------------------------------------------------

\subsubsection{Regulatory Frameworks and their Assumptions}
\label{sec:gov_regulatory}

The FDA's authorized AI device landscape \citep{warraich2024fda} and the EU
AI Act both operate from a shared premise: that governance is a precondition
that health systems and AI developers must satisfy before market entry.
Under this model, authorization requires demonstrating safety and efficacy
through premarket evaluation; post-market obligations are limited largely to
adverse event reporting. \citet{tian2026pilottrap} challenge this premise
empirically: their 18-month implementation study reveals that the governance
capacity required for sustainable AI deployment is not, in fact, pre-existing
in health systems. Accountability boundaries, coordination mechanisms, and
infrastructure governance structures emerged through implementation — they
were discovered, not pre-installed.

This finding has a direct implication for regulatory design: a governance
framework that treats capacity as a precondition rather than an outcome of
the implementation process will systematically misidentify the point of
intervention. Regulators specifying what governance should achieve without
attending to whether health systems have the organizational infrastructure
to deliver it will produce well-intentioned but structurally ineffective
governance standards.

% -----------------------------------------------------------------------------
% 2.2 Pilot trap: organizational barriers to sustainable deployment
% -----------------------------------------------------------------------------

\subsubsection{The Pilot Trap as an Organizational Failure}
\label{sec:gov_pilottrap}

The pilot trap — defined as the failure of clinical AI applications to
transition from short-term demonstration projects to sustained components
of routine clinical care — appears across multiple health system contexts
and AI application types. Tian et al. \citep{tian2026pilottrap} attribute
this not to algorithmic underperformance but to three organizational
deficits: (\textit{i}) unclear accountability boundaries (who is legally
and professionally responsible for an AI-assisted clinical decision?);
(\textit{ii}) absent coordination mechanisms (how do clinical, technical,
and administrative units interact around an AI system?); and
(\textit{iii}) inadequate infrastructure governance (how are data pipelines,
model updates, and system integrations managed operationally?).

These organizational deficits are not unique to a single health system or
national context. \citet{keri2026simulated} documents analogous challenges
in psychiatric AI deployment, where the risks of simulated empathy, illusion
of attachment, and unpredictable crisis-mode behavior by AI chatbots require
governance mechanisms — including strict validation protocols, continuous
monitoring, and clear human accountability — that most psychiatric services
are not currently equipped to provide.

% -----------------------------------------------------------------------------
% 2.3 Adaptive regulatory approaches
% -----------------------------------------------------------------------------

\subsubsection{Emerging Adaptive and Participatory Governance Models}
\label{sec:gov_adaptive}

In response to the limitations of static pre-market regulatory models,
several papers in the reviewed corpus describe adaptive governance
approaches that embed learning and iteration into the regulatory process
itself. Regulatory sandboxes — already adopted in fintech and now proposed
for health AI — allow AI systems to be deployed in controlled real-world
settings with enhanced monitoring, time-limited authorization, and
mandatory performance thresholds that must be met for full market entry
\citep{pham2025ethical, abdul2025ai}. The WHO's ethics and governance
guidance on AI in health similarly advocates for continuous compliance
mechanisms that evolve with the AI system's performance history rather
than treating authorization as a one-time gatekeeping event.

Participatory governance mechanisms — incorporating affected communities,
patient advocacy groups, and frontline clinical staff into governance
processes — are proposed across multiple reviewed papers as essential
for AI legitimacy and sustainability \citep{bouderhem2025shaping,
pham2025ethical}. This aligns with the broader literature on democratic
AI and public interest technology, which holds that the epistemic
diversity brought by participatory processes improves both the quality
and the legitimacy of governance decisions.

In response to the limitations of static pre-market regulatory models,
several papers describe adaptive governance approaches that embed learning
and iteration into the regulatory process itself. Regulatory sandboxes —
already adopted in fintech and now proposed for health AI — allow AI systems
to be deployed in controlled real-world settings with enhanced monitoring,
time-limited authorization, and mandatory performance thresholds that must
be met for full market entry \citep{pham2025ethical, abdul2025ai}. The FDA's
Predetermined Change Control Plan (PCCP) represents a parallel mechanism,
permitting manufacturers to define in advance the types of model changes that
can be implemented without requiring new market authorization, subject to
documented performance monitoring. The EU AI Act's conformity assessment
procedure similarly incorporates post-market surveillance obligations and
mandatory periodic reviews, though the governance capacity to operationalize
these requirements remains underdeveloped in practice. The WHO's ethics and
governance guidance on AI in health advocates for continuous compliance
mechanisms that evolve with the AI system's performance history rather than
treating authorization as a one-time gatekeeping event.

% =============================================================================
% THEME 3: IMPLEMENTATION OUTCOMES — Barriers and Societal Impact
% =============================================================================

\subsection{Implementation Outcomes: Barriers and Societal Impact}
\label{sec:results_implementation}

% -----------------------------------------------------------------------------
% 3.1 Documented barriers to sustainable implementation
% -----------------------------------------------------------------------------

\subsubsection{Organizational, Technical, and Social Barriers}
\label{sec:impl_barriers}

The implementation barriers documented in the reviewed literature can be
organized into three analytically distinct but interacting categories.

\textbf{Organizational barriers} include: workforce trust deficits (clinicians
skeptical of AI recommendations that conflict with their clinical judgment);
misaligned incentive structures (AI systems optimized for metrics that do
not reflect clinical priorities or institutional goals); change management
failures (inadequate training, poor human-AI workflow integration); and
institutional governance deficits (absence of clear ownership, accountability,
or escalation protocols for AI-assisted decisions)
\citep{tian2026pilottrap, pham2025ethical}.

\textbf{Technical barriers} include: data quality and interoperability
problems (AI systems trained on curated research data that does not reflect
the noise, incompleteness, and heterogeneity of real Electronic Health Record
[EHR] data); infrastructure limitations (insufficient computing resources,
network reliability, or system integration capacity in resource-constrained
settings); and model update challenges (managing the regulatory, technical,
and organizational complexity of updating a deployed AI system)
\citep{warraich2024fda, abulibdeh2025illusion}.

\textbf{Social barriers} include: algorithmic literacy gaps among clinicians
and patients; power dynamics that privilege technical AI teams over
frontline clinical expertise; and the risk of deskilling — where reliance
on AI systems erodes the development of clinical skills that are needed
when AI is unavailable or fails.

% -----------------------------------------------------------------------------
% 3.2 Equity and bias concerns
% -----------------------------------------------------------------------------

\subsubsection{Algorithmic Bias, Equity, and Differential Performance}
\label{sec:impl_equity}

The equity dimension of health AI implementation receives significant
attention across the reviewed corpus. \citet{abulibdeh2025illusion} document
that AI-enabled medical devices authorized by the FDA have historically
underrepresented minority populations in their training and validation data,
creating a systematic risk that AI systems perform differentially across
racial, socioeconomic, and geographic subgroups. This bias amplification
risk is particularly acute in clinical domains such as dermatology, radiology,
and cardiovascular risk prediction, where training data skews toward
Caucasian and male populations. The absence of mandatory subgroup performance
reporting in current regulatory frameworks means that differential AI
performance may go undetected for years.

\citet{pham2025ethical} argue that addressing algorithmic bias requires not
only technical interventions (diverse data collection, bias auditing) but
also governance mechanisms that institutionalize equity considerations
throughout the AI lifecycle — from problem framing and dataset construction
to deployment monitoring and post-market evaluation.

% -----------------------------------------------------------------------------
% 3.3 Evidence quality and research gaps
% -----------------------------------------------------------------------------

\subsubsection{Evidence Quality and Research Gaps}
\label{sec:impl_gaps}

A consistent finding across the reviewed literature is the limited
high-quality evidence on real-world clinical impact of AI systems. The
majority of published studies report technical performance metrics
(accuracy, AUC) in controlled settings. Prospective, multicenter studies
measuring patient outcomes (mortality, morbidity, length of stay, quality
of life) attributable to AI-assisted care remain rare. This evidence gap
is compounded by publication bias: AI systems that perform well in
validation studies are more likely to be published, while failed
deployments are rarely reported, creating a systematically distorted
literature that overstates real-world effectiveness.

The three thematic findings are analytically connected rather than parallel.
The ecological validity gap does not merely reflect a technical measurement
problem — it has direct consequences for governance. When AI systems are
validated under conditions that do not reflect real clinical complexity,
governance frameworks based on those validations will be structurally
misaligned with the phenomena they seek to govern. Similarly, the governance
capacity deficit is not simply a matter of underinvestment — it is partly a
consequence of evaluation frameworks that generate overconfident validation
evidence, reducing the perceived urgency of building pre-deployment governance
infrastructure. The pilot trap can then be understood as the downstream
consequence of these upstream failures: systems validated on inadequate
evidence, governed by frameworks that assume absent capacity, predictably
fail to transition to routine care. This interdependency implies that
improving health AI outcomes requires simultaneous action across all three
dimensions — better evaluation methodology, governance capacity investment,
and implementation support — rather than sequential approaches that treat
evaluation, governance, and implementation as separate phases.